% =========================================================
%  SECCIÓN II — MARCO TEÓRICO
% =========================================================
\section{Marco Te\'orico}

\subsection{El Operador Prewitt y la Magnitud del Gradiente}

El operador Prewitt~\cite{prewitt1970} es un operador diferencial
discreto que aproxima el gradiente de la funci\'on de intensidad
de una imagen en escala de grises.
A diferencia del operador Sobel, el Prewitt utiliza pesos uniformes
(todos iguales a~1) en sus m\'ascaras de convoluci\'on, lo que lo
convierte en el m\'as simple de la familia de operadores basados
en gradiente.

Las dos m\'ascaras del operador son:

\begin{equation}
  K_x = \begin{pmatrix}
    -1 & 0 & 1 \\
    -1 & 0 & 1 \\
    -1 & 0 & 1
  \end{pmatrix}, \quad
  K_y = \begin{pmatrix}
    +1 & +1 & +1 \\
     0 &  0 &  0 \\
    -1 & -1 & -1
  \end{pmatrix}
  \label{eq:mascaras}
\end{equation}

donde $K_x$ detecta cambios en la direcci\'on horizontal (bordes
verticales) y $K_y$ detecta cambios en la direcci\'on vertical
(bordes horizontales).
Tras calcular los gradientes $G_x$ e $G_y$ por convoluci\'on de
la imagen con las m\'ascaras, la magnitud del gradiente en cada
p\'ixel se obtiene como:

\begin{equation}
  G[i] = \sqrt{G_x[i]^2 + G_y[i]^2}
  \label{eq:magnitud}
\end{equation}

Esta f\'ormula constituye el \textit{kernel num\'erico} objeto de
este estudio.
Para cada elemento $i$ del arreglo de datos se requieren exactamente
cuatro operaciones de punto flotante: dos multiplicaciones
($G_x^2$, $G_y^2$), una adici\'on y una ra\'iz cuadrada.
Por tanto, la intensidad aritm\'etica te\'orica del kernel,
entendida como el cociente entre operaciones de punto flotante
(\textit{FLOPs}) y transferencias de memoria (\textit{Bytes}), es:

\begin{equation}
  \mathcal{I} = \frac{4 \text{ FLOPs}}{3 \times 4 \text{ Bytes}}
             = \frac{4}{12} \approx 0.333 \text{ FLOP/Byte}
  \label{eq:ai}
\end{equation}

donde el denominador refleja la lectura de los dos arreglos de
entrada ($G_x$, $G_y$) y la escritura del arreglo de salida
($\textit{Mag}$), cada uno de tipo \texttt{float} (4~Bytes).

\subsection{Vectorizaci\'on SIMD}

La vectorizaci\'on SIMD es una t\'ecnica de paralelismo de datos
que permite ejecutar la misma operaci\'on aritm\'etica sobre
m\'ultiples datos simult\'aneamente utilizando registros vectoriales
anchos~\cite{reinders2007}.
En los procesadores Intel modernos existen tres extensiones SIMD
principales relevantes para este trabajo:

\begin{itemize}
  \item \textbf{SSE/SSE2} (\textit{Streaming SIMD Extensions}):
    registros XMM de 128~bits, procesa hasta 4~floats por
    instrucci\'on.
  \item \textbf{AVX/AVX2} (\textit{Advanced Vector Extensions}):
    registros YMM de 256~bits, procesa hasta 8~floats por
    instrucci\'on.
  \item \textbf{AVX-512}: registros ZMM de 512~bits, procesa
    hasta 16~floats por instrucci\'on.
\end{itemize}

La vectorizaci\'on puede obtenerse mediante tres estrategias,
en orden creciente de control por parte del programador:

\begin{enumerate}
  \item \textbf{Autom\'atica}: el compilador analiza el c\'odigo
    escalar y genera instrucciones SIMD cuando es seguro hacerlo.
    Requiere indicios como el calificador \texttt{\_\_restrict}
    para descartar aliasing de punteros.

  \item \textbf{Guiada}: el programador a\~nade pragmas o
    directivas que instruyen al compilador a vectorizar un bucle
    determinado con un ancho espec\'ifico.
    La directiva \texttt{\#pragma omp simd simdlen(8)} indica
    al compilador que genere SIMD con 8 elementos por vector.

  \item \textbf{Expl\'icita}: el programador escribe directamente
    las instrucciones vectoriales mediante intr\'insecos
    (\textit{intrinsics}), funciones C que se mapean uno a uno con
    instrucciones ensamblador SIMD.
    Por ejemplo, \texttt{\_mm256\_sqrt\_ps} calcula la ra\'iz
    cuadrada de 8~floats en paralelo usando el registro YMM.
\end{enumerate}

La principal ventaja te\'orica de la vectorizaci\'on es un
incremento del rendimiento proporcional al ancho del vector:
$8\times$ con AVX2 respecto al c\'omputo escalar.
Sin embargo, en la pr\'actica este factor est\'a limitado por
el ancho de banda de memoria del sistema~\cite{sutter2005}.

\subsection{Intel Advisor y el Modelo Roofline}

Intel Advisor~\cite{intel_advisor} es una herramienta de an\'alisis
de rendimiento orientada a la optimizaci\'on de vectorizaci\'on
y paralelismo.
Ofrece dos tipos principales de an\'alisis:

\begin{itemize}
  \item \textbf{Survey}: identifica los bucles del programa,
    determina si est\'an vectorizados o no, reporta el tiempo de
    ejecuci\'on y el conteo de iteraciones (\textit{trip count}).

  \item \textbf{Tripcounts + FLOP}: mide el n\'umero real de
    operaciones de punto flotante ejecutadas y el volumen de
    datos transferidos desde/hacia memoria, calculando as\'i la
    intensidad aritm\'etica medida del programa.
\end{itemize}

El \textbf{modelo Roofline}~\cite{williams2009} es un marco
visual para analizar el rendimiento de un kernel en funci\'on
de su intensidad aritm\'etica.
Define los l\'imites te\'oricos alcanzables como:

\begin{equation}
  P_{\max} = \min\!\left(P_{\text{peak}},\;
             \mathcal{I} \times B_{\text{mem}}\right)
  \label{eq:roofline}
\end{equation}

donde $P_{\max}$ es el rendimiento m\'aximo (GFLOPS),
$P_{\text{peak}}$ es el pico computacional del procesador,
$\mathcal{I}$ es la intensidad aritm\'etica (FLOP/Byte) y
$B_{\text{mem}}$ es el ancho de banda de memoria (GB/s).

Un kernel es \textit{memory-bound} cuando su punto de operaci\'on
cae a la izquierda del ``punto de cresta'' (\textit{ridge point}),
es decir, cuando la memoria no puede suministrar datos a la
velocidad que el procesador puede calcular.
En ese caso, incrementar las instrucciones SIMD no mejora el
rendimiento, ya que el cuello de botella est\'a en el bus de
datos y no en la unidad aritm\'etica.
